\paragraph{Complete finite-resolution energy readout.}
For a unit one-particle probe $f$, the same ground defines the probability
measure on signed addition and removal energies
\[
 \mu_f(B)=\langle Vf,1_B(L)Vf\rangle.
\]
The complete Kraus pair $c(f)^\dagger,c(f)$ labels addition
and removal; each branch has a nonnegative physical excitation energy.
Only the response-frequency variable assigns a minus sign to removal.
The exact identities $LV=Vh+UJ$ and
$V^*J=J^*J=rI$ give the operator moments
\[
 M_0=I,\qquad M_1=h+UrI,\qquad
 M_2=h^2+2Urh+U^2rI.
\]
The second-moment identity is a quadratic-form statement and does not
require $Vf\in\operatorname{Dom}L^2$.
For a local mode at $c=1/8,U=1$, the mean lies in $[1/8,1/2]$, the second
moment in $[19/128,67/128]$, and the variance in $[17/128,35/128]$.
Thus $\mu_f(|E|\ge R)\le67/(128R^2)$.
More generally, $\|hf\|\le K$ gives a second moment at most
$K^2+UK+U^2/2$, useful for uniformly controlled physical smearing.

Choose $M\ge2$ and $\tau>0$.  The positive effects
\[
 K_k(E)=\left|\frac1M\sum_{j=0}^{M-1}
 e^{-ij(\tau E-2\pi k/M)}\right|^2,\qquad
 q_k=\int K_k(E)\,d\mu_f(E),\quad 0\le k<M,
\]
sum exactly to one by finite Fourier orthogonality.  Writing
$\phi_f(t)=\int e^{-itE}\,d\mu_f(E)$, their complete probabilities are
\[
 q_k=\frac1M+\frac2M\sum_{n=1}^{M-1}
 \left(1-\frac nM\right)\operatorname{Re}
 \bigl(e^{2\pi ikn/M}\phi_f(n\tau)\bigr).
\]
The principal energy centers are modulo $2\pi/\tau$ and are readout bins,
not asserted eigenvalues.  For $0<\tau\Delta\le3$, circular error has
probability at most $\min\{1,10/(M\tau^2\Delta^2)\}$, improved to
$22/(7M\tau\Delta-44)$ when its denominator is positive.  This follows
by summing the two inverse-square grid tails.  A second-moment bound $C$
adds an alias term at most $C/(3/\tau-\Delta)^2$.  Consequently the
probability of a real energy error at least $\Delta$ is controlled,
without a pointwise spectral density or an isolated-pole assumption.

The complete frozen 78-line cube measure gives certified distributions
for $M=8,16,64$ at $\tau=1$.  At $\Delta=1/2$, direct positive-kernel
and Fourier interval calculations bound the respective error probabilities
by $0.100284517428$, $0.055193158232$, and $0.020209031434$.
All outcomes are retained.  The earlier executed native time-one
quadratures check the first harmonic; higher harmonics here are source
calculations, not new native long-time or finite-shot measurements.

Along the existing full-ground subsequence, the finite positive spectral
measures converge weakly.  The uniform fixed-spacing spatial kernel bound
also promotes entrywise convergence to strong convergence of the
compressed resolvents on the common one-particle space.  No identification
of the varying GNS Hilbert spaces is assumed.  An explicit finite readout
error adds the local-dynamics, native-control, and statistical errors to
the unspecified finite-ground state discrepancy; the last term is
never set to zero.  Sharp interval probabilities require zero limiting
mass at their endpoints.  These results do not supply physical particle
identifications, relativistic rest masses, or an interacting continuum.
